% !PRINTED 437
\weberpage{437}
%<<CONT
\[
2^{\frac{n-1}{2}} \prod\limits_{1,\,\frac{n-1}{2}}^{\nu} \cos\frac{2\nu m\pi}{n} = \pm 1, \tag{19}
\]
in which for the time being the sign still remains undetermined.

But if in (18) we let $\varphi$ pass into $\pi:2$, the right-hand side, which presents itself under the indeterminate form $0:0$, acquires the limiting value $(-1)^{\frac{n-1}{2}} n$, and we obtain again, generally for every odd $n$, the formula which we derived in the preceding section for a prime $n$
\[
2^{\frac{n-1}{2}} \prod\limits_{1,\,\frac{n-1}{2}}^{\nu} \sin\frac{2\nu m\pi}{n} = \pm \sqrt{n}, \tag{20}
\]
and for $m=1$ the positive sign is to be taken.

The function $B_n$, which with respect to $x$ is of degree $(n-1)$, vanishes when $n\varphi$, but not $\varphi$, is a multiple of $\pi$, that is, for
\[
x = 2\cos\frac{h\pi}{n}, \tag{21}
\]
in which $h$ is an integer not divisible by $n$. If we first assume $n$ even, we obtain all the \gesperrt{distinct} values contained in (21) by taking for $m$ an arbitrary number coprime to $n$ and setting
\[
h = m,\ 2m,\ 3m \ldots (n-1)m \tag{22}
\]
for
\[
\cos\frac{h\pi}{n} = \cos\frac{h'\pi}{n},
\]
holds only when either $h+h'$ or $h-h'$ is a multiple of $2n$, which cannot occur when $h$ and $h'$ are both taken from the sequence (22).

Among the values (21) there is one, namely the one corresponding to the value $h = nm:2$, which is equal to zero, and the rest are opposite in pairs. Accordingly, for an \gesperrt{even} $n$ we obtain
\[
B_n(x) = x \prod\limits_{1,\,\frac{n}{2}-1}^{\nu} \left(x^2 - 4\cos^2\frac{\nu m\pi}{n}\right). \tag{23}
\]

If $n$ is odd, then $B_n(x)$ depends only on $x^2$; hence the values (21) are equal and opposite in pairs (for $h$ and $h+n$). If, again supposing $m$ relatively prime to $n$, one
%CONT>>
